\section{Two transverse labels per matrix edge}

\subsection*{Finding O-027: the exact determinant row threshold when \(\lambda_3\leq2\)}
Let \(n\geq2\). For every outer degree \(d\) and every
\(\lambda\vdash nd\) with
\[
                    \ell(\lambda)\leq2n,\qquad \lambda_3\leq2,
                                                               \tag{O-027.1}
\]
the determinant orbit closure retains all ambient copies:
\[
                         K_n(\lambda)=a_\lambda(d[n]).
                                                               \tag{O-027.2}
\]
For the unpadded permanent of size \(n\), the same conditions give
\[
                         q^n_\lambda(d[n])=a_\lambda(d[n]).
                                                               \tag{O-027.3}
\]
These statements include all highest-weight multiplicity copies.
They follow from a single explicit construction which makes the complete
projection to transverse order two along a binary plane dominant for
both determinant and permanent in \(2n\) variables.
The construction improves O-020's sufficient determinant size
\(n\geq r-1\) to \(n\geq\lceil r/2\rceil\).
O-020 and its integral incidence calculation remain valid independently.

For \(n\geq3\), the row threshold \(2n\) in the universal assertion
(O-027.2) is sharp. The known LMR module has
\[
 \Lambda_n=\bigl(2n^3-4n^2+1,\ 2n^2-4n+1,\ 2^{\,2n-1}\bigr),
 \qquad d=2n(n-1),
                                                               \tag{O-027.4}
\]
so that \(\ell(\Lambda_n)=2n+1\) and \((\Lambda_n)_3=2\).
Theorem 2.3.1 and Section 3.2 of Landsberg--Manivel--Ressayre
give a nonzero copy of this module in the determinant ideal.
Consequently \(K_n(\Lambda_n)<a_{\Lambda_n}(d[n])\).
This sharpness statement concerns ambient determinant multiplicity;
it does not assert a permanent comparison for all \(n\).
The finite permanent endpoint certificates in A-031 supply that
additional comparison for their stated sizes.

\paragraph{The binary pencil and its quadratic coefficient space.}
Put \(N=n-1\), and use the variables
\[
                 x,y,\ z_{i1},z_{i2}\quad(1\leq i\leq N).
                                                               \tag{O-027.5}
\]
Let \(V=\bigoplus_{i=1}^{N}V_i\), where
\(V_i=\langle z_{i1},z_{i2}\rangle\), and choose distinct scalars
\(a_1,\ldots,a_n\). Write
\[
                \ell_i=y-a_i x,\qquad \Delta=\prod_{i=1}^n\ell_i.
\]
The desired coefficient space is
\[
 \mathcal J=
 \operatorname{Sym}^n\langle x,y\rangle
 \ \oplus\
 \operatorname{Sym}^{n-1}\langle x,y\rangle\otimes V
 \ \oplus\
 \operatorname{Sym}^{n-2}\langle x,y\rangle\otimes\operatorname{Sym}^2V.
                                                               \tag{O-027.6}
\]
Projection to \(\mathcal J\) discards only terms of total \(z\)-degree
at least three. Higher discarded coefficients are unrestricted.

For \(i<j<n\), set
\[
 A^0_{ij}=z_{i1}+2z_{j1},\qquad
 A^0_{ji}=z_{i2}+2z_{j2}.
\]
For \(i<n\), set
\[
                  A^0_{in}=z_{i1},\qquad A^0_{ni}=z_{i2},
\]
and set \(A^0_{ii}=\ell_i\).
Thus \(A^0\) is a completely specified matrix of linear forms.
The quadratic transverse part of its permanent is
\[
       \sum_{i<j}\frac{\Delta}{\ell_i\ell_j}A^0_{ij}A^0_{ji};
                                                               \tag{O-027.7}
\]
the corresponding determinant part is its negative.
Indeed, a permutation with exactly two off-diagonal entries is a
transposition. All cycles of length at least three disappear in
the projection (O-027.6).

\paragraph{The two-label incidence lemma.}
Let
\[
 H_n=\{(c_1,\ldots,c_n):\textstyle\sum_i c_i=0\}.
\]
For \(i<j<n\), put
\[
 U_{ij}=\langle z_{i1}+2z_{j1},\,z_{i2}+2z_{j2}\rangle,
 \qquad U_{in}=V_i.
\]
The linear map
\[
 \Psi:\bigoplus_{i<j}(U_{ij}\otimes V)
       \longrightarrow H_n\otimes\operatorname{Sym}^2V,
 \qquad
 (u\otimes v)_{ij}\longmapsto(e_i-e_j)\otimes uv
                                                               \tag{O-027.8}
\]
is surjective over every field of characteristic different from
two and three.

\begin{proof}
A functional on the target can be represented by symmetric bilinear
forms \(Q_1,\ldots,Q_n\) on \(V\), modulo adding the same form to every
\(Q_i\). Set \(Q_n=0\). A functional annihilates the image of \(\Psi\)
exactly when
\[
             (Q_i-Q_j)(u,v)=0
       \quad\text{for all }u\in U_{ij},\ v\in V.
                                                               \tag{O-027.9}
\]
The reference edges \(\{i,n\}\) first imply that \(Q_i\) kills \(V_i\)
in either argument.

For \(i<j<n\), apply (O-027.9) to
\(u=z_i+2z_j\), where the coordinates of \(z_i\in V_i\) and
\(z_j\in V_j\) agree under the displayed bases.
Taking the other argument in \(V_i\) shows that the
\(V_i\)-diagonal block of \(Q_j\) is zero.
Taking it in \(V_j\), and dividing by two, shows that the
\(V_j\)-diagonal block of \(Q_i\) is zero.
Together with the reference edges, this kills every diagonal
two-by-two block of every \(Q_i\).

Every block still possibly nonzero is indexed by three distinct
nonreference vertices. Fix \(i<j<k<n\), and denote by
\[
            A=(Q_i)_{jk},\qquad B=(Q_j)_{ik},\qquad C=(Q_k)_{ij}
\]
their two-by-two matrices. The edges \(\{i,j\}\), \(\{i,k\}\),
and \(\{j,k\}\), respectively, give
\[
                         B=2A,\qquad C=2A^{\mathsf T},
                         \qquad C=2B.
                                                               \tag{O-027.10}
\]
Hence \(A^{\mathsf T}=2A\), and transposition gives \(A=4A\).
Characteristic different from three forces \(A=0\), and then
\(B=C=0\). This treats every remaining block.
There is no triple to treat when \(N<3\); the preceding diagonal
argument already gives the same conclusion.
The dual kernel is zero, proving surjectivity.
\end{proof}

\paragraph{From binary residues to the full coefficient differential.}
The map
\[
 \operatorname{Sym}^{n-2}\langle x,y\rangle\longrightarrow H_n,
 \qquad
 f\longmapsto
 \left(\frac{f(1,a_i)}{\prod_{j\ne i}(a_i-a_j)}\right)_{i=1}^n
                                                               \tag{O-027.11}
\]
is an isomorphism. Its image is in \(H_n\) by the coefficient
of \(y^{-1}\) in the partial-fraction expansion of
\(f(1,y)/\Delta(1,y)\). The map is injective because a polynomial
of degree at most \(n-2\) vanishing at all \(n\) distinct \(a_i\)
is zero; source and target both have dimension \(n-1\).
In particular,
\[
 \frac{\Delta}{\ell_i\ell_j}\longmapsto
                    \frac{e_i-e_j}{a_i-a_j}.
                                                               \tag{O-027.12}
\]

Vary a single off-diagonal entry \(A_{ij}\) by an arbitrary
linear form \(H\in V\). The transverse-order-two part of the
derivative is exactly
\[
              \frac{\Delta}{\ell_i\ell_j}\,A^0_{ji}H
                                                               \tag{O-027.13}
\]
for the permanent and its negative for the determinant.
It has no transverse degree zero or one. Varying both
orientations of every edge, using (O-027.12), gives precisely
the image in (O-027.8), up to nonzero scalar factors.
The incidence lemma therefore supplies the whole third
summand of (O-027.6).

Next vary \(A_{ii}\) by \(H\in V\). Up to transverse order two,
the derivative is
\[
                         \frac{\Delta}{\ell_i}H.
                                                               \tag{O-027.14}
\]
There is no transverse-degree-one term in the remaining
principal minor, since a nonidentity permutation cannot have
exactly one off-diagonal entry. The \(n\) binary forms
\(\Delta/\ell_i\) form a basis in degree \(n-1\), as evaluation
at \(a_i\) shows. These variations supply the entire middle
summand, without affecting the other two.

Finally vary each diagonal entry by \(x\) or \(y\).
The transverse-degree-zero images span
\(\operatorname{Sym}^n\langle x,y\rangle\), because they are
\(x\Delta/\ell_i\) and \(y\Delta/\ell_i\).
Their possible transverse-degree-two components can be
subtracted using (O-027.13); their transverse-degree-one
components are zero. The complete differential onto
\(\mathcal J\) is consequently surjective, for both determinant
and permanent.

\paragraph{Dominance and all-copy consequence.}
The coefficients of a matrix of linear forms are affine
parameters, and projection of its determinant or permanent
to (O-027.6) is a polynomial map.
A surjective differential at \(A^0\) gives an injective
coordinate pullback: restrict source directions to a linear
right inverse of the differential, translate source and
target to their base values, and use the lowest nonzero
homogeneous part of a polynomial under a map
\(t\mapsto t+O(t^2)\).
Thus the coefficient projection is dominant.
Setting unnecessary variables to zero gives the same
conclusion for every \(r\leq2n\): coefficients using only
the retained variables are independent of the discarded
matrix-entry coefficients, so variable restriction commutes
with projection to the retained target coordinates.

For any \(\lambda\) satisfying (O-027.1), every highest-weight
vector in the degree-\(d\) ambient coordinate ring uses only
the first \(\ell(\lambda)\) variable directions.
The flag-minor support theorem F-016 says that every coefficient
generator appearing in such a vector has total exponent at
most \(\lambda_3\leq2\) in coordinates three and later.
It is therefore a polynomial in the coordinates of
\(\mathcal J\).
If the vector lies in the determinant ideal, it vanishes
on every matrix-pencil substitution into the determinant,
since such substitutions belong to the orbit closure.
The injective pullback just proved forces it to be zero.
The same reasoning applies to the unpadded permanent.
Every vector in the full highest-weight ideal multiplicity
space is covered by this argument. Consequently that space
is zero, proving (O-027.2)--(O-027.3).

\paragraph{Internal padding.}
The argument above varies all matrix-entry coefficients.
In particular, it does not claim that \(A_{11}=x\) can be
fixed in the same construction.
For the internally padded permanent, the verified O-021
padding argument instead gives the following consequence:
if \(n\geq m\), \(\ell(\lambda)\leq r\), \(\lambda_3\leq2\),
and
\[
                  m\geq
                \max\{\lceil r/2\rceil+1,\ \lambda_2+1\},
                                                               \tag{O-027.15}
\]
then
\[
                           q^m_\lambda(d[n])
                         =K_n(\lambda)=a_\lambda(d[n]).
                                                               \tag{O-027.16}
\]
Indeed, put \(s=m-1\). The new construction supplies the
full permanent coefficient projection in size \(s\), because
\(r\leq2s\). The other inequality gives \(s\geq\lambda_2\),
so F-014 stabilizes the complete highest-weight space at
base degree \(s\). The diagonal, invertible differential
transpose in O-012 preserves its transverse-order-two
coefficient support. Finally the block matrix
\(\operatorname{diag}(x,A_s)\) has fixed internal entry \(x\)
and permanent \(x\operatorname{per}(A_s)\), exactly as in
O-021. This proves all-copy equality after padding.
For one-row partitions, the nonzero pure-power coordinate
gives the same equality directly.

\paragraph{Sharp boundary for the determinant coefficient projection.}
For \(r\) variables, the target in transverse order two has dimension
\[
 D_n(r)=(n+1)+n(r-2)+(n-1)\binom{r-1}{2}.
                                                               \tag{O-027.17}
\]
The closure of the \(r\)-variable determinant-pencil image
has dimension at most \((r-2)n^2+2\), by O-010's
left-right gauge argument. At \(r=2n+1\), subtraction gives
\[
             D_n(2n+1)-\bigl((2n-1)n^2+2\bigr)=n-1>0.
                                                               \tag{O-027.18}
\]
Hence the full coefficient projection cannot be dominant
there. For larger \(r\), the difference increases.
Together with the construction, this makes \(r=2n\) the
exact maximal number of variables for this determinant
coefficient projection, within the ambient variable range.
This is a geometric statement about the whole projection;
the separate representation-theoretic sharpness in
(O-027.4) uses the known LMR ideal copy.

\paragraph{The selected integral incidence cokernel.}
The incidence map (O-027.8) is defined over \(\mathbb Z\).
With \(N=n-1\), its finite cokernel is
\[
 \operatorname{coker}\Psi_{\mathbb Z}
 \simeq
 (\mathbb Z/2)^{\,3\binom N2+3\binom N3}
       \ \oplus\
 (\mathbb Z/6)^{\,\binom N3}.
                                                               \tag{O-027.19}
\]
Here binomial coefficients are zero outside their usual range.
To see this, use the dual matrix solely to compute Smith
invariants, which are unchanged by transposition.
The reference-edge equations eliminate, by unit pivots,
every block of \(Q_i\) involving \(V_i\).
For each pair \(i<j\), the \(V_i\)-diagonal block of \(Q_j\)
is eliminated by unit pivots, while the three independent
entries of the \(V_j\)-diagonal block of \(Q_i\) have the
relations \(2q=0\).
For each triple \(i<j<k\), use the unit relations
\(B=2A\) and \(C=2B\) in (O-027.10).
The remaining four-entry block is
\[
                          2(A^{\mathsf T}-2A)=0.
\]
Its two diagonal entries contribute \(2,2\); on its two
off-diagonal entries the matrix is
\[
                    \begin{pmatrix}-4&2\\2&-4\end{pmatrix},
\]
whose Smith invariants are \(2,6\).
All pairs and triples involve disjoint remaining variables,
giving (O-027.19).
This is the cokernel of the selected incidence map;
it is not a claim that every full determinant or permanent
jet map must fail in characteristic two or three.

\paragraph{Grothendieck dictionary and base-change scope.}
The matrix pencil, determinant, permanent, coefficient
projection and incidence homomorphism are polynomial or
linear constructions over \(\mathbb Z\).
Equation (O-027.19) makes the integral limitation explicit.
After inverting six and the differences \(a_i-a_j\),
the incidence and binary interpolation maps are split
surjective, and the full differential has a linear
right inverse. The resulting lowest-term coordinate
injection survives arbitrary commutative base change
over this localized ring.
The passage to irreducible highest-weight multiplicities
is made over \(\mathbb C\), using F-016 and
characteristic-zero semisimplicity.
In Grothendieck's language, the graded coordinate algebras
and their ideals are comodules as in SGA~3, I,
Propositions~4.7.2--4.7.3, pp.~34--36.
Their Schur operations belong to the special
\(\lambda\)-structures of SGA~6, VI, Theorem~3.3,
pp.~392--393. The conclusion is certified by the actual
coordinate injection and its explicit integral localization,
not by positivity of a virtual \(K\)-class.
